\lesson{30}{Mar 18 2022 Fri (23:06:17)}{Computing expected value in a game of chance}{Unit 7}

\begin{formula}[Expected Value Formula]
  \label{frm:expected_value_formula}

  \begin{equation}
    \label{eqt:expected_value_formula}
    \begin{split}
      E = x_1 \times p_1 + x_2 \times p_2 + \ldots + x_n \times p_n
    \end{split}
  .\end{equation}

  The values $x_1, x_2, \ldots, x_n$ are all of the possible results.

  The values $p_1, p_2, \ldots, p_n$ are the corresponding probabilities of the
  results.
\end{formula}

Each time Reuben decides to play the game of chance, the result is uncertain.
However, the expected value can be used to predict what could happen if Reuben
plays the game many times.

To find the expected value of a game of chance, we multiply each possible result
(in dollars) by its probability and then add these products.

\begin{mdframed}
  \begin{itemize}
    \label{item:questions}

    \item[(a)] There are $8$ possible results of playing the game. Each result
      is a dollar amount. Since the dart will land on a numbered slice at
      random, each result is equally likely to occur. So, the probability is
      $\frac{1}{8}$.

      \begin{table}[H]
        \centering
        \caption{Expected Value}
        \label{tab:expected_value_in_table_form}

        \begin{tabular}{|l|c|c|c|c|c|c|c|c|}
          \hline
          Dart Toss & $1$ & $2$ & $3$ & $4$ & $5$ & $6$ & $7$ & $8$ \\
          \hline
          Result & $\$1$ & $\$2$ & $\$5$ & $\$8$ & $\$-4$ & $\$-4$ & $\$-4$ &
            $\$-4$ \\
          \hline
          Probability & $\frac{1}{8}$ & $\frac{1}{8}$ & $\frac{1}{8}$ &
            $\frac{1}{8}$ & $\frac{1}{8}$ & $\frac{1}{8}$ & $\frac{1}{8}$ &
            $\frac{1}{8}$ \\
          \hline
        \end{tabular}
      \end{table}

      Using the formula for expected value, we get the following.

      \begin{equation}
        \label{eqt:expected_value_sequence}
        \begin{alignedat}{3}
          E  = &1 \left(\frac{1}{8}\right) +
                2 \left(\frac{1}{8}\right) +
                5 \left(\frac{1}{8}\right) +
                8 \left(\frac{1}{8}\right) - \\
               &4 \left(\frac{1}{8}\right) -
                4 \left(\frac{1}{8}\right) -
                4 \left(\frac{1}{8}\right) -
                4 \left(\frac{1}{8}\right) \\
            &= \frac{1}{8} + \frac{2}{8} + \frac{5}{8} \frac{8}{8} - \frac{4}{8}
              - \frac{4}{8} - \frac{4}{8} - \frac{4}{8} \\
            &= \frac{1 + 2 + 5 + 8 - 4 - 4 - 4 - 4}{8} \\
            &= \$0
        \end{alignedat}
      .\end{equation}

      So, the expected value of playing the game is $\$0$.

      Notice that it's impossible to get a result of $\$0$ any particular time
      the game is played. This is just an indicator of the outcome after
      multiple games have been played.

    \item[(b)] 	What can Reuben expect in the long run, after playing the game
      many times?

      The expected value is zero. So, Reuben can expect to break even (neither
      gain nor lose money).
  \end{itemize}
\end{mdframed}

\newpage
